\subsection{Fibrations}

6.1.1. A fibration of class \(C^{r}\) , or simply fibration, is a triple \((X, B, \pi)\) where B and X are varieties of class \(C^{r}\) and \(\pi\) is a morphism from X to B, having the following property:

\begin{enumerate}
\def\labelenumi{(\Alph{enumi})}
\setcounter{enumi}{5}
\tightlist
\item
  For every \(x \in \mathbf{B}\) , there exists an open neighborhood \(\mathbf{U}\) of \(x\) , a manifold \(\mathbf{F}\) , and an isomorphism \(\varphi\) of \(\pi^{-1}(\mathbf{U})\) onto \(\mathbf{U} \times \mathbf{F}\) such that \(\pi(\varphi^{-1}(x, y)) = x\) for all \(x \in \mathbf{U}\) and every \(y \in \mathbf{F}\) .
\end{enumerate}

If \(\lambda = (X, B, \pi)\) is a fibration, one calls \(X\) the space of \(\lambda\) , \(B\) the base of \(\lambda\) , and \(\pi\) the projection of \(\lambda\) . The map \(\pi\) is a submersion; in particular \(\pi(X)\) is open in \(B\) and if \(R\) denotes the equivalence relation defined by \(\pi\) on \(X\) , the canonical map from \(X/R\) to \(\pi(X)\) is an isomorphism. For every \(x \in B\) the inverse image \(\pi^{-1}(x)\) is a closed submanifold of \(X\) , called the fiber of \(x\) , and denoted \(X_x\) .

6.1.2. Examples:

\begin{enumerate}
\def\labelenumi{(\alph{enumi})}
\item
  If B and F are two manifolds, the triple \((\mathbf{B} \times \mathbf{F}, \mathbf{B}, \mathbf{pr}_{1})\) is a fibration whose fibers are canonically isomorphic to F.
\item
  If \(\lambda = (\mathbf{X},\mathbf{B},\pi)\) and \(\lambda^{\prime} = (\mathbf{X}^{\prime},\mathbf{B}^{\prime},\pi^{\prime})\) are fibrations,
\end{enumerate}

\[
\lambda \times \lambda^ {\prime} = (\mathrm{X} \times \mathrm{X} ^ {\prime}, \mathrm{B} \times \mathrm{B} ^ {\prime}, \pi \times \pi^ {\prime})
\]

is a fibration; it is called the product of \(\lambda\) and of \(\lambda'\) .

\begin{enumerate}
\def\labelenumi{(\alph{enumi})}
\setcounter{enumi}{2}
\tightlist
\item
  If \(\lambda = (X, B, \pi)\) and \(\lambda' = (X', B, \pi')\) are fibrations with the same base, \(\lambda \times_{B} \lambda' = (X \times_{B} X', B, \pi \times_{B} \pi')\) is a fibration; it is called the product of \(\lambda\) and of \(\lambda'\) above \(B\) , or also the fibered product of \(\lambda\) and of \(\lambda'\) .
\end{enumerate}

6.1.3. Let \(\lambda = (X, B, \pi)\) and \(\lambda' = (X', B', \pi')\) be two fibrations. One calls a morphism of \(\lambda\) to \(\lambda'\) any pair \((f, g)\) , where \(f\) is a morphism of \(B\) to \(B'\) , and \(g\) a morphism of \(X\) to \(X'\) such that \(\pi' \circ g = f \circ \pi\) . When \(B = B'\) and \(f = \mathrm{Id}_B\) , one says that \(g\) is a \(B\) -morphism of \(\lambda\) to \(\lambda'\) ; if \(g\) is an isomorphism of \(X\) onto \(X'\) , \(g^{-1}\) is a \(B\) -morphism of \(\lambda'\) to \(\lambda\) , and one says that \(g\) is a \(B\) -isomorphism from \(\lambda\) onto \(\lambda'\) ; for this to be the case, it is necessary and sufficient that, for every \(x \in B\) , the map \(g_x: X_x \to X'_x\) induced by \(g\) be an isomorphism.

6.1.4. The fibration \((\mathbf{B} \times \mathbf{F}, \mathbf{B}, \mathbf{pr}_{1})\) is called the trivial fibration with base B and fiber F. An isomorphism of a fibration \(\lambda\) onto a trivial fibration is called a trivialization of \(\lambda\) .

6.1.5. Let \(\lambda = (X, B, \pi)\) be a fibration, and let \(f: B' \to B\) be a morphism. Let \(\pi'\) be the canonical morphism of \(B' \times_B X\) to \(B'\) . The triple \((B' \times_B X, B', \pi')\) is a fibration, called the inverse image of \(\lambda\) by \(f\) or the fibration deduced from \(\lambda\) by base change of \(B\) to \(B'\) along \(f\) , and denoted \(B' \times_B \lambda\) , or again \(f^*\lambda\) . If \(f'\) denotes the canonical map of \(B' \times_B X\) to \(X\) , the pair \((f, f')\) is a morphism of \(B' \times_B \lambda\) to \(\lambda\) ; it enjoys the following universal property: if \((f, g)\) is a morphism of a fibration \(\lambda'\) with base \(B'\) into the fibration \(\lambda\) , there exists a unique \(B'\) -morphism \(\varphi: \lambda' \to B' \times_B \lambda\) such that \((f, g) = (f, f') \circ \varphi\) .

When B' is a submanifold of B, and f is the canonical injection of B' into B, \(B' \times_{B} X\) is identified with the submanifold \(\pi^{-1}(B')\) of X, and \(\pi'\) is identified with the restriction of \(\pi\) to \(\pi^{-1}(B')\) ; the inverse image of \(\lambda\) by f is then called the fibration induced by \(\lambda\) on B'.

6.1.6. If \(\lambda = (X, B, \pi)\) is a fibration, one calls a morphic section (or simply section) of \(\lambda\) every morphism \(s: B \to X\) such that \(\pi \circ s = \mathrm{Id}_B\) .

